\documentclass[aspectratio=169,10pt]{beamer}

% ------------------------------------------------
% pdfLaTeX: lietuviški simboliai ir šriftai
% ------------------------------------------------
\usepackage[T1]{fontenc}
\usepackage[utf8]{inputenc}
\usepackage[lithuanian,english]{babel}
\usepackage[sfdefault]{noto}

% Matematika
\usepackage{amsmath,amssymb,cancel}

% ------------------------------------------------
% Beamer tema ir spalvos
% ------------------------------------------------
\usetheme{Madrid}
\usecolortheme{whale}

\setbeamertemplate{navigation symbols}{}

% Lietuviški Beamer aplinkų pavadinimai
\deftranslation[to=lithuanian]{Definition}{Apibrėžimas}
\deftranslation[to=lithuanian]{Example}{Pavyzdys}
\deftranslation[to=lithuanian]{Remark}{Pastaba}
\deftranslation[to=lithuanian]{Proof}{Įrodymas}

% Papildomos aplinkos ir komandos
\newenvironment{property}[1][]%
{\begin{block}{Savybė: #1}}%
{\end{block}}

\newcommand{\vocab}[1]{\textbf{\color{structure.fg}#1}}

% ------------------------------------------------
% Dokumento informacija
% ------------------------------------------------
\title[Laipsniai su natūraliuoju rodikliu]
{Laipsnis su natūraliuoju rodikliu ir jo savybės}

\subtitle{Teorija, įrodymai, pavyzdžiai ir uždaviniai}

\author[A. Berniukevičius]{Andrius Berniukevičius}

\date{\today}

% =================================================
\begin{document}
% =================================================


% ------------------------------------------------
% Titulinė skaidrė
% ------------------------------------------------
\begin{frame}
    \titlepage
\end{frame}


% ------------------------------------------------
% Laipsnio apibrėžimas
% ------------------------------------------------
\begin{frame}{Laipsnio apibrėžimas}

    \begin{definition}[Laipsnis su natūraliuoju rodikliu]

        Skaičiaus $a$ \vocab{laipsniu su natūraliuoju rodikliu $n$}
        (kai $n>1$) vadinama $n$ vienodų dauginamųjų,
        kurių kiekvienas lygus $a$, sandauga:

        \[
        a^n
        =
        \underbrace{
        a\cdot a\cdot a\cdot\dots\cdot a
        }_{n\text{ kartų}}
        \]

    \end{definition}

    \vspace{0.2cm}

    \begin{block}{Svarbios sąvokos ir taisyklės}

        \begin{itemize}

            \item Pirmojo laipsnio taisyklė:
            \[
            a^1=a
            \]

            \item
            \textbf{$a$} – \vocab{laipsnio pagrindas}
            (skaičius, kurį dauginame iš savęs).

            \item
            \textbf{$n$} – \vocab{laipsnio rodiklis}
            (rodo, kiek kartų pagrindas kartojasi).

        \end{itemize}

    \end{block}

\end{frame}


% ------------------------------------------------
% Daugyba ir kėlimas laipsniu
% ------------------------------------------------
\begin{frame}{Daugyba vs. Kėlimas laipsniu}

    \begin{alertblock}{Svarbi pastaba}
        Būtina atskirti skaičių daugybą nuo kėlimo laipsniu:
    \end{alertblock}

    \vspace{0.3cm}

    \begin{columns}[T]

        \begin{column}{0.48\textwidth}

            \begin{block}{Daugyba (Sudedame kartus)}

                \[
                4\cdot3=4+4+4=12
                \]

                Skaičių $4$ \textbf{sudedame} $3$ kartus.

            \end{block}

        \end{column}

        \begin{column}{0.48\textwidth}

            \begin{block}{Kėlimas laipsniu (Sudauginame kartus)}

                \[
                4^3=4\cdot4\cdot4=64
                \]

                Skaičių $4$ \textbf{sudauginame} iš savęs $3$ kartus.

            \end{block}

        \end{column}

    \end{columns}

\end{frame}


% ------------------------------------------------
% 1 savybė
% ------------------------------------------------
\begin{frame}{1 Savybė: Laipsnių daugyba}

    \begin{property}[Vienodų pagrindų daugyba]

        Kai dauginami laipsniai su vienodais pagrindais,
        pagrindas paliekamas tas pats, o rodikliai sudedami:

        \[
        a^m\cdot a^n=a^{m+n}
        \]

    \end{property}

    \begin{proof}

        Užrašome abu sandaugos narius pagal apibrėžimą:

        \[
        a^m\cdot a^n
        =
        \left(
        \underbrace{
        a\cdot a\cdot\dots\cdot a
        }_{m\text{ kartų}}
        \right)
        \cdot
        \left(
        \underbrace{
        a\cdot a\cdot\dots\cdot a
        }_{n\text{ kartų}}
        \right)
        \]

        Kadangi skliaustai daugyboje esmės nekeičia,
        iš viso skaičius $a$ pasikartoja $m+n$ kartų:

        \[
        a^m\cdot a^n
        =
        \underbrace{
        a\cdot a\cdot\dots\cdot a
        }_{m+n\text{ kartų}}
        =
        a^{m+n}
        \]

    \end{proof}

\end{frame}


% ------------------------------------------------
% 2 savybė
% ------------------------------------------------
\begin{frame}{2 Savybė: Laipsnių dalyba}

    \begin{property}[Vienodų pagrindų dalyba]

        Kai dalijami laipsniai su vienodais pagrindais,
        pagrindas paliekamas tas pats, o iš dalinio rodiklio
        atimamas daliklio rodiklis:

        \[
        a^m:a^n=a^{m-n}
        \qquad
        (\text{kai }a\neq0\text{ ir }m>n)
        \]

    \end{property}

    \begin{proof}

        Užrašome dalybą trupmena ir suprastiname
        $n$ vienodų dauginamųjų skaitiklyje ir vardiklyje:

        \[
        \frac{a^m}{a^n}
        =
        \frac{
        \overbrace{
        a\cdot a\cdot\dots\cdot a
        }^{m\text{ kartų}}
        }{
        \underbrace{
        a\cdot a\cdot\dots\cdot a
        }_{n\text{ kartų}}
        }
        =
        \underbrace{
        a\cdot a\cdot\dots\cdot a
        }_{m-n\text{ kartų}}
        =
        a^{m-n}
        \]

    \end{proof}

\end{frame}


% ------------------------------------------------
% Dalybos pavyzdys
% ------------------------------------------------
\begin{frame}{Dalybos savybės pavyzdys}

    \begin{example}[Kompaktiškas suprastinimas]

        Apskaičiuokime $\dfrac{a^7}{a^3}$:

        \[
        \frac{a^7}{a^3}
        =
        \frac{
        \cancelto{1}{a}\cdot
        \cancelto{1}{a}\cdot
        \cancelto{1}{a}\cdot
        a\cdot a\cdot a\cdot a
        }{
        \cancelto{1}{a}\cdot
        \cancelto{1}{a}\cdot
        \cancelto{1}{a}
        }
        =
        \frac{
        1\cdot1\cdot1\cdot
        a\cdot a\cdot a\cdot a
        }{1}
        =
        a^4
        \]

    \end{example}

    \begin{block}{Išvada}

        Iš pradinių $7$ dauginamųjų skaitiklyje
        atėmėme $3$ dauginamuosius iš vardiklio,
        todėl galutinis laipsnio rodiklis gaunamas atimties būdu:

        \[
        a^{7-3}=a^4
        \]

    \end{block}

\end{frame}


% ------------------------------------------------
% 3 savybė
% ------------------------------------------------
\begin{frame}{3 Savybė: Laipsnio kėlimas laipsniu}

    \begin{property}[Laipsnio kėlimas laipsniu]

        Keliant laipsnį laipsniu, pagrindas paliekamas tas pats,
        o rodikliai sudauginami:

        \[
        (a^m)^n=a^{m\cdot n}
        \]

    \end{property}

    \begin{proof}

        Pagal apibrėžimą pagrindą $(a^m)$ dauginame iš savęs
        $n$ kartų:

        \[
        (a^m)^n
        =
        \underbrace{
        a^m\cdot a^m\cdot\dots\cdot a^m
        }_{n\text{ kartų}}
        \]

        Pritaikome laipsnių daugybos savybę
        (rodiklius sudedame):

        \[
        (a^m)^n
        =
        a^{
        \overbrace{
        m+m+\dots+m
        }^{n\text{ kartų}}
        }
        =
        a^{m\cdot n}
        \]

    \end{proof}

\end{frame}


% ------------------------------------------------
% 4 savybė
% ------------------------------------------------
\begin{frame}{4 Savybė: Sandaugos kėlimas laipsniu}

    \begin{property}[Sandaugos kėlimas laipsniu]

        Keliant sandaugą laipsniu, kiekvienas dauginamasis
        keliamas tuo laipsniu atskirai:

        \[
        (a\cdot b)^n=a^n\cdot b^n
        \]

    \end{property}

    \begin{proof}

        Sugrupuojame dauginamuosius $a$ ir $b$ atskirai:

        \[
        (a\cdot b)^n
        =
        \underbrace{
        (a\cdot b)\cdot\dots\cdot(a\cdot b)
        }_{n\text{ kartų}}
        =
        \left(
        \underbrace{
        a\cdot a\cdot\dots\cdot a
        }_{n\text{ kartų}}
        \right)
        \cdot
        \left(
        \underbrace{
        b\cdot b\cdot\dots\cdot b
        }_{n\text{ kartų}}
        \right)
        =
        a^n\cdot b^n
        \]

    \end{proof}

\end{frame}


% ------------------------------------------------
% 5 savybė
% ------------------------------------------------
\begin{frame}{5 Savybė: Trupmenos kėlimas laipsniu}

    \begin{property}[Trupmenos kėlimas laipsniu]

        Keliant trupmeną laipsniu, skaitiklis ir vardiklis
        keliami tuo laipsniu atskirai:

        \[
        \left(\frac{a}{b}\right)^n
        =
        \frac{a^n}{b^n}
        \qquad (b\neq0)
        \]

    \end{property}

    \begin{proof}

        Pagal trupmenų daugybos taisyklę:

        \[
        \left(\frac{a}{b}\right)^n
        =
        \underbrace{
        \frac{a}{b}\cdot
        \frac{a}{b}\cdot
        \dots\cdot
        \frac{a}{b}
        }_{n\text{ kartų}}
        =
        \frac{
        \overbrace{
        a\cdot a\cdot\dots\cdot a
        }^{n\text{ kartų}}
        }{
        \underbrace{
        b\cdot b\cdot\dots\cdot b
        }_{n\text{ kartų}}
        }
        =
        \frac{a^n}{b^n}
        \]

    \end{proof}

\end{frame}


% ------------------------------------------------
% Pavyzdys 1
% ------------------------------------------------
\begin{frame}{Pavyzdys 1: Reiškinio supaprastinimas}

    \begin{example}

        Supaprastinkite reiškinį:

        \[
        \left(x^{45}:x^{25}\right)^2
        \cdot
        \left(x^5\right)^3
        :
        x^{10}
        \]

    \end{example}

    \begin{block}{Sprendimo žingsniai}

        \begin{enumerate}

            \item
            Veiksmas skliaustuose:
            \[
            x^{45}:x^{25}
            =
            x^{45-25}
            =
            x^{20}
            \]

            \item
            Kėlimas laipsniu:
            \[
            (x^{20})^2
            =
            x^{20\cdot2}
            =
            x^{40}
            \]

            \item
            Antrasis narys:
            \[
            (x^5)^3
            =
            x^{5\cdot3}
            =
            x^{15}
            \]

            \item
            Daugyba ir dalyba:
            \[
            x^{40}\cdot x^{15}:x^{10}
            =
            x^{40+15-10}
            =
            x^{45}
            \]

        \end{enumerate}

    \end{block}

    \begin{alertblock}{Atsakymas}
        \[
        x^{45}
        \]
    \end{alertblock}

\end{frame}


% ------------------------------------------------
% Pavyzdys 2
% ------------------------------------------------
\begin{frame}{Pavyzdys 2: Skaitinis reiškinys}

    \begin{example}

        Apskaičiuokite reiškinio reikšmę:

        \[
        \frac{
        \left(
        \frac{243^5\cdot9^8}{81^3\cdot27^6}
        \right)^4
        \cdot
        \left(3^{10}\right)^2
        }{
        \left(
        \frac{729^2}{9^4}
        \right)^3
        \cdot
        3^{12}
        }
        \]

    \end{example}

    \begin{block}{Sprendimas (visus laipsnius pertvarkome į pagrindą $3$)}
        \small
        \begin{columns}[T,onlytextwidth]
            \begin{column}{0.48\textwidth}
                \vspace{0.4em}
                \textbf{Skaitiklis:}
                \[
                    \begin{aligned}
                    \left(\frac{3^{25}\cdot3^{16}}{3^{12}\cdot3^{18}}\right)^4
                    \cdot3^{20}
                    &=(3^{11})^4\cdot3^{20}\\
                    &=3^{64}
                    \end{aligned}
                \]
            \end{column}
            \begin{column}{0.48\textwidth}
                \vspace{0.4em}
                \textbf{Vardiklis:}
                \[
                    \begin{aligned}
                    \left(\frac{3^{12}}{3^8}\right)^3\cdot3^{12}
                    &=(3^4)^3\cdot3^{12}\\
                    &=3^{24}
                    \end{aligned}
                \]
                \textbf{Galutinė dalyba:} $3^{64}:3^{24}=3^{40}$
            \end{column}
        \end{columns}
    \end{block}

\end{frame}


% ------------------------------------------------
% Uždaviniai 1–13
% ------------------------------------------------
\begin{frame}{Uždaviniai (1 dalis: 1--13)}

    Atlikite veiksmus su laipsniais:

    \vspace{0.2cm}

    \begin{columns}[T]

        \begin{column}{0.5\textwidth}

            \begin{enumerate}

                \item
                $(((x^2)^5)^4)^3:x^{100}$

                \item
                $\left(\frac{a^{25}}{b^{15}}\right)^4
                \cdot
                \left(\frac{b^{30}}{a^{40}}\right)^2$

                \item
                $(y^{15}\cdot y^{25})^3:(y^{12})^8\cdot y^{10}$

                \item
                $(a^{10}b^{20}c^{30})^5:
                (a^{12}b^{24}c^{36})^4$

                \item
                $\frac{
                (p^{18}:p^{10})^5\cdot p^{20}
                }{
                (p^3\cdot p^7)^4
                }$

                \item
                $\frac{
                (2a^{15})^4\cdot(5a^{10})^2
                }{
                (10a^{30})^2
                }$

                \item
                $\left(
                \frac{x^{25}\cdot y^{35}}
                {x^{15}\cdot y^{20}}
                \right)^4$

            \end{enumerate}

        \end{column}

        \begin{column}{0.5\textwidth}

            \begin{enumerate}

                \setcounter{enumi}{7}

                \item
                $\frac{u^{50}}{v^{30}}
                :
                \left(
                \frac{u^{15}}{v^{10}}
                \right)^3$

                \item
                $x^{100}
                :
                \left(
                \frac{
                (x^{15})^4\cdot x^{20}
                }{
                x^{10}
                }
                \right)$

                \item
                $\frac{
                (m^{12}n^{18})^5\cdot k^{40}
                }{
                (m^{20}n^{30})^3\cdot(k^8)^5
                }$

                \item
                $(((z^3)^4)^2\cdot z^{15}):z^{30}$

                \item
                $\left(
                \frac{c^{40}}{d^{20}}
                \right)^3
                :
                \left(
                \frac{c^{50}}{d^{30}}
                \right)^2$

                \item
                $\frac{
                (3x^{12}y^8)^4
                }{
                (9x^{20}y^{15})^2
                }$

            \end{enumerate}

        \end{column}

    \end{columns}

\end{frame}


% ------------------------------------------------
% Uždaviniai 14–26
% ------------------------------------------------
\begin{frame}{Uždaviniai (2 dalis: 14--26)}

    \begin{columns}[T]

        \begin{column}{0.5\textwidth}

            \begin{enumerate}

                \setcounter{enumi}{13}

                \item
                $\frac{
                (a^{15}\cdot a^{25})^4
                }{
                (a^{20})^7\cdot a^{10}
                }$

                \item
                $(w^{50}:w^{20})^3
                \cdot
                (w^{10}:w^5)^4$

                \item
                $\left(
                \frac{
                x^{30}y^{40}z^{50}
                }{
                x^{20}y^{25}z^{45}
                }
                \right)^5$

                \item
                $\frac{
                (2^{30}\cdot2^{40})^2
                }{
                (2^{35})^4
                }$

                \item
                $(a^{18}\cdot b^{25})^5:
                (a^{15}\cdot b^{20})^6$

                \item
                $\frac{
                (((p^5)^4)^3)^2
                }{
                p^{100}\cdot p^{15}
                }$

                \item
                $\left(
                \frac{m^{60}}{n^{40}}
                \right)^2
                :
                \frac{
                m^{100}
                }{
                (n^{16})^5
                }$

            \end{enumerate}

        \end{column}

        \begin{column}{0.5\textwidth}

            \begin{enumerate}

                \setcounter{enumi}{20}

                \item
                $\frac{
                (4^5\cdot8^3)^2
                }{
                16^6\cdot2^{10}
                }$

                \item
                $\frac{
                (9^{10}\cdot27^5)^2
                }{
                (81^4)^3\cdot3^{15}
                }$

                \item
                $\frac{
                125^8\cdot25^{10}
                }{
                (5^6)^4\cdot25^5
                }$

                \item
                $\left(
                \frac{
                32^4\cdot8^6
                }{
                4^{10}\cdot2^{15}
                }
                \right)^3$

                \item
                $\frac{
                (1000^5\cdot100^8)^2
                }{
                10^{50}\cdot10000^2
                }$

                \item
                $\frac{
                64^5\cdot(16^3)^4
                }{
                (4^8)^3\cdot8^5
                }$

            \end{enumerate}

        \end{column}

    \end{columns}

\end{frame}


% ------------------------------------------------
% Atsakymai
% ------------------------------------------------
\begin{frame}{Atsakymai (1--26)}

    \begin{columns}[T]

        \begin{column}{0.25\textwidth}
            \small

            (1) $x^{20}$\\
            (2) $a^{20}$\\
            (3) $y^{34}$\\
            (4) $a^2b^4c^6$\\
            (5) $p^{20}$\\
            (6) $4a^{20}$\\
            (7) $x^{40}y^{60}$

        \end{column}

        \begin{column}{0.25\textwidth}
            \small

            (8) $u^5$\\
            (9) $x^{30}$\\
            (10) $1$\\
            (11) $z^9$\\
            (12) $c^{20}$\\
            (13) $x^8y^2$\\
            (14) $a^{10}$

        \end{column}

        \begin{column}{0.25\textwidth}
            \small

            (15) $w^{110}$\\
            (16) $x^{50}y^{75}z^{25}$\\
            (17) $1$\\
            (18) $b^5$\\
            (19) $p^5$\\
            (20) $m^{20}$\\
            (21) $16$

        \end{column}

        \begin{column}{0.25\textwidth}
            \small

            (22) $2187$\\
            (23) $5^{10}$\\
            (24) $512$\\
            (25) $10000$\\
            (26) $32768$

        \end{column}

    \end{columns}

\end{frame}


% =================================================
\end{document}